\chapter{What a Computer Can Do}\label{ch:turing}

\begin{etymology}
\textbf{Compute} comes from the Latin \emph{computare}: \emph{com-},
``together'', and \emph{putare}, ``to reckon'' (and, originally, ``to prune'',
to clean something up in order to assess it). To compute was to count by
putting things together; the word only later came to mean processing
information by rules.

\textbf{Complexity} comes from \emph{complexus}: \emph{com-} plus
\emph{plectere}, ``to braid''. Something complex is braided together and hard to
pull apart. In computing it means the cost of solving a problem and how that
cost grows with the size of the input.
\end{etymology}

Saying a quantum computer is ``faster'' means nothing until you say faster than
what, at what, and measured how. That needs a precise model of what a classical
computer is. The model was built in 1936, ten years before the first electronic
computer, to answer a question about the foundations of mathematics.

\section{The Turing machine}

\begin{history}{1928--1936 \textperiodcentered{} Göttingen, Cambridge and Princeton}
In 1928, David Hilbert and Wilhelm Ackermann posed the
\emph{Entscheidungsproblem}, the decision problem: is there a mechanical
procedure that decides, for any statement of first-order logic, whether it is
provable? To answer ``no'', someone had to define what a mechanical procedure
is. Alonzo Church, at Princeton, did it with the lambda calculus in 1936. Alan
Turing, a 23-year-old fellow at King's College, Cambridge, did it the same year
with an imaginary machine, in ``On Computable Numbers''. Both answers were no.
\end{history}

Turing's machine has four parts: an infinite tape divided into cells, a head that
reads and writes one cell at a time, a finite set of internal states, and a
transition table. At each step the table looks at the current state and the
symbol under the head and says three things: which symbol to write, which way to
move, and which state to go to next.

A machine that inverts every bit of a binary string needs one working state and
three rules: on a 0, write 1 and move right; on a 1, write 0 and move right; on a
blank, stop. Fed \texttt{110100}, it walks across the tape and leaves
\texttt{001011}. Nothing about this looks like a computer, and that is the point:
anything a laptop does can be broken down into steps this small.

\begin{concept}{Church--Turing thesis}
Anything that can be computed by a mechanical procedure can be computed by a
Turing machine.

\textbf{What it fixes:} it gives ``computable'' a precise meaning, so you can
prove that some things are not.

\textbf{What it costs:} it is a thesis, not a theorem. It cannot be proved,
because ``mechanical procedure'' is an informal idea. It has held up for ninety
years.
\end{concept}

Turing also showed that one machine can imitate any other if you write the other
machine's table on its tape. That \emph{universal} machine is the idea behind
software: a program is the description of a machine, and the computer is the
universal machine that runs it.

\section{What no machine can do}

Turing's real target was a limit. He proved that no machine can decide, for every
machine $M$ and input $x$, whether $M$ eventually stops on $x$. This is the
\emph{halting problem}. The proof is a diagonal argument, like Cantor's: suppose
such a decider exists, build a machine that asks the decider about itself and
then does the opposite, and you get a contradiction.

So there are well-posed questions that no algorithm can answer, ever. A quantum
computer does not change this. It computes exactly the same functions as a Turing
machine; what it might change is how fast.

\section{Cost: P, NP and their cousins}

Computable is not the same as practical. A problem can have an algorithm that
would take longer than the age of the universe. Complexity theory sorts problems
by how the cost grows with the size $n$ of the input.

\begin{history}{1965--1972 \textperiodcentered{} the birth of P and NP}
In 1965, Jack Edmonds and Alan Cobham, independently, proposed that ``efficient''
should mean ``polynomial time'': the number of steps grows like $n^k$ for some
fixed $k$. In 1971, Stephen Cook showed that a large family of problems that are
easy to check are all equally hard to solve, if any of them is. Richard Karp
listed 21 of them the next year. P versus NP was born, and it is still the most
famous open problem in computer science.
\end{history}

The classes that matter here:

\begin{itemize}
\item \textbf{P}: problems solvable in polynomial time by a deterministic machine.
\item \textbf{NP}: problems where a proposed solution can be checked in polynomial
time. Factoring is in NP: given $p$ and $q$, multiply and compare.
\item \textbf{BPP}: problems solvable in polynomial time by a machine that may flip
coins, with error probability below $1/3$ (repeating the run makes it as small as
you like). Defined by John Gill in 1977.
\item \textbf{BQP}: the same thing for a quantum computer. Defined by Bernstein and
Vazirani in 1993.
\end{itemize}

It is proved that $\mathrm{P} \subseteq \mathrm{BPP} \subseteq \mathrm{BQP}$:
anything a classical machine does efficiently, a quantum machine does too. What
nobody has proved is that any of these inclusions is strict. Factoring is in BQP
thanks to Shor, and nobody knows whether it is in P. It is also not believed to
be NP-complete, which matters: Shor's algorithm says nothing about the hard
problems in Karp's list.

The gap between polynomial and exponential is the whole story. At $n = 50$, $n^2$
is 2500 steps and $2^n$ is about $1.1 \times 10^{15}$.

\section{The quantum Turing machine}

Deutsch's 1985 machine keeps Turing's parts but lets the tape, the head and the
state be in superposition, and replaces the transition table with a unitary
transformation (Chapter~\ref{ch:state}). Nobody programs with it. In 1989
Deutsch also introduced quantum circuits, sequences of gates acting on qubits,
and in 1993 Andrew Yao proved that the two models are equivalent. Circuits are
what everybody uses.

The machine can be in a superposition of many configurations, but a measurement
returns only one of them. The advantage is never ``try everything and read
everything''. It is to arrange the computation so that wrong answers cancel each
other before the measurement and the right one remains. Chapter~\ref{ch:qubit}
shows what that means with one qubit.

\section{What the classes say about quantum speedups}

Shor put factoring in BQP. Grover gave a quadratic speedup for search, and in
1997 Bennett, Bernstein, Brassard and Vazirani proved that quadratic is the best
possible for unstructured search. A quantum computer cannot solve NP-complete
problems by brute force in polynomial time.

So the exponential speedups come from structure, like the periodicity Shor
exploits. Where there is no structure, you get Grover's square root at best.
Whether BQP is really bigger than BPP, outside of oracle models, is still open.
Shor's algorithm is the strongest evidence that it is, but it is evidence, not
proof: nobody has shown that factoring is classically hard.
